% Generated draft: requires mathematical and source review.
\documentclass[11pt]{article}
\usepackage[margin=25mm]{geometry}
\usepackage{fontspec}
\setmainfont{Latin Modern Roman}
\usepackage{amsmath,amssymb,mathtools,unicode-math}
\setmathfont{Latin Modern Math}
\usepackage{enumitem,needspace,xurl,hyperref}
\setlength{\parindent}{0pt}
\setlength{\parskip}{5pt}
\setlength{\emergencystretch}{3em}
\newcommand{\N}{\mathbb{N}}
\newcommand{\Z}{\mathbb{Z}}
\newcommand{\Q}{\mathbb{Q}}
\newcommand{\R}{\mathbb{R}}
\newcommand{\C}{\mathbb{C}}
\newcommand{\F}{\mathbb{F}}
\newcommand{\A}{\mathbb{A}}
\newcommand{\PP}{\mathbb{P}}
\newcommand{\E}{\mathbb{E}}
\newcommand{\eps}{\varepsilon}
\newcommand{\dd}{\,\mathrm{d}}
\begin{document}
\section*{UnsolvedMath ID: 30006521}
\textbf{Source problem code:} catalogue entry
\subsection{Definitions and mathematical statement}
For each positive integer \(n\), let \(S_n\) be the symmetric group on \(\{1,\ldots,n\}\). Embed these groups by adjoining fixed points, and write \(S_\infty=\bigcup_{n\geq1}S_n\). For a permutation \(u\), let \(\ell(u)\) be its inversion number, and let \(s_k=(k\;k+1)\). Define strict left weak order by
\[
u<_L v
\quad\Longleftrightarrow\quad
v=s_{i_r}\cdots s_{i_1}u\text{ for some }r\geq1,
\qquad \ell(v)=\ell(u)+r.
\]
Under the customary interpretation in terms of ordinary single Grothendieck polynomials and their expansion in Schubert polynomials, write
\[
G_w=\sum_{\sigma\in S_\infty}c_{w,\sigma}\,\mathfrak S_\sigma,
\qquad
\operatorname{Supp}_S(G_w)=\{\sigma\in S_\infty:c_{w,\sigma}\neq0\},
\]
where \(\mathfrak S_\sigma\) denotes a Schubert polynomial and the sum has finite support. Let \(\deg(G_w)\) denote total polynomial degree.

With these conventions, the cleaned question in the supplied record is whether, for every positive integer \(n\), every \(w\in S_n\), and every \(u\in\operatorname{Supp}_S(G_w)\),
\[
\ell(u)<\deg(G_w)
\quad\Longrightarrow\quad
\exists v\in\operatorname{Supp}_S(G_w)\text{ such that }u<_L v.
\]
The supplied record does not explicitly define the Grothendieck-polynomial convention, Schubert support, or ambient permutation index set. The interpretation above is therefore conditional and requires checking against the report; in particular, it does not impose an unsupported requirement that all support permutations belong to \(S_n\).

The original excerpt also presents a proposed stronger formulation, labeled Conjecture~2: under the same hypothesis, there exists \(k\) such that \(u<s_ku\) and \(s_ku\in\operatorname{Supp}_S(G_w)\). The relation \(<\) and the range of \(k\) are not specified in the supplied excerpt. If the intended condition is a left weak cover, its precise form, with the ambient \(S_\infty\) convention above, is
\[
\exists k\geq1:\qquad
\ell(s_ku)=\ell(u)+1,
\qquad s_ku\in\operatorname{Supp}_S(G_w).
\]
This second formulation must remain distinguished from the cleaned target, which only asks for a strictly larger support element. Its precise interpretation and the excerpt's conjecture numbering require source review.
\subsection{Short English statement}
The question asks whether every permutation in the Schubert support of \(G_w\) whose length is below the total degree has a strictly larger support element in left weak order. The original excerpt also asks for an extension by one left simple reflection. Polynomial conventions, the ambient support domain, and the second formulation's order notation require verification.
\subsection{Sources}
\begin{itemize}
\item Supplied local problem registry: permanent UnsolvedMath ID \(30006521\), title ``Upper Elements in Grothendieck Schubert Supports,'' exact source problem code \texttt{catalogue entry}. The registry labels the statement \texttt{corrected\_verified}, qualified as repairs only to formula delimiters, notation, grammar, and typography. No individual proposer is identified; the supplied proposed year is 2026.
\item \emph{Enumerative Combinatorics} (2026), Oberwolfach Reports, pp.~81--156, DOI \texttt{10.4171/owr/2026/2}: \url{https://doi.org/10.4171/owr/2026/2}. The supplied background describes these reports as records of mathematical workshops and their open-problem discussions.
\item Mészáros--Setiabrata--St.~Dizier, \emph{On the support of Grothendieck polynomials}: \url{https://arxiv.org/abs/2201.09452}.
\item Chou--Yu, \emph{Grothendieck polynomials of inverse fireworks permutations}: \url{https://arxiv.org/abs/2403.12186}.
\item Supplied literature assessment, checked at the exact recorded date \texttt{2026-08-22}: status ``partially solved.'' The assessment states that the stronger left-weak-order support statement is established only in special permutation classes and that its general form remains open. This is supplied dated curation derived from the OWR record, not an independently verified assessment. The supplied background expressly calls for independent verification of time-sensitive status and cited claims.
\end{itemize}
\end{document}
